\documentclass[11pt]{article}

\usepackage[margin=1in]{geometry}
\usepackage{amsmath,amssymb,amsthm,mathtools}
\usepackage{enumitem}
\usepackage{hyperref}
\hypersetup{colorlinks=true,linkcolor=blue,citecolor=blue,urlcolor=blue}
\usepackage{orcidlink}
\usepackage{fontawesome5}

\newtheorem{claim}{Claim}[section]

\newcommand{\ZZ}{\mathbb{Z}}
\newcommand{\NN}{\mathbb{N}}

\title{Errors in Claimed Complete Proofs (and Disproofs)\\
of the Collatz Conjecture: A Catalog and Taxonomy\\
of Twelve Cases}

\author{Renato Augusto Tavares~{\small\href{mailto:dr.renatotavares@gmail.com}{\textcolor{black}{\faEnvelope}}}~\orcidlink{0009-0002-0196-3311}\\
\normalsize Universidade Federal de Goiás}
\date{\today}

\begin{document}

\maketitle

\begin{abstract}
The Collatz conjecture attracts an unusually large volume of claimed
complete proofs and disproofs, mostly circulated through preprint
servers and low-scrutiny venues rather than peer-reviewed mathematics
journals. We report a systematic review of twelve such claims
encountered during a broader literature survey, each read in full,
independently reconstructed, and checked computationally wherever the
claim admits a numerical test. None of the twelve is a valid proof or
disproof. The contribution of this note is not the count itself but a
taxonomy of the ways they fail: two categories account for ten of the
twelve cases (circular reasoning that smuggles the conjecture's
conclusion into a definition or premise, three cases; illegitimate
generalization from one hand-picked adversarial sequence to all
sequences, two cases), and a single recurring fallacy --- treating an
aggregate, average, or asymptotic property as a deterministic,
pointwise guarantee for one specific trajectory --- accounts for five
cases by itself, disguised in turn by Baker's theorem, a random-walk
heuristic, a 2-adic Lyapunov argument, Big-O notation, and static
2-adic density. The remaining two cases are an elementary error
unrelated to the conjecture's real difficulty, and a genuine rigor gap
that, unlike its closest analogue in the catalog, produced no
computational counterexample. Ten of twelve failures converge, by
superficially unrelated routes, on the same structural obstruction that
no known method resolves: the impossibility of bounding adversarial
runs of increasing steps without, at some point, assuming what is to be
proved.
\end{abstract}

\section{Introduction}
\label{sec:intro}

The Collatz conjecture has been open for close to ninety years, and
its statement is simple enough that it draws an unusually large number
of claimed complete solutions from outside the community of
professional number theorists. Almost all of them are circulated
through preprint repositories, personal websites, or journals with no
established reputation in pure mathematics, rather than peer-reviewed
venues with specialist referees. None has been accepted by the
mathematical community. This is not, by itself, evidence that any
particular claim is wrong --- but it sets the correct prior: a claimed
complete proof (or disproof) of a problem this resistant, from a venue
without specialist peer review, should be read with the working
assumption that it contains a genuine gap, most likely at exactly the
point where the known ``structural obstruction'' to the conjecture
would have to be overcome (the impossibility, with any method known
today, of bounding how long an adversarial run of increasing steps can
be --- see \cite{LagariasSurvey} for the standard account of why this
is hard).

This note catalogs twelve such claims, encountered over the course of
a broader literature survey conducted alongside an unrelated research
program on a $qx+1$ generalization of Tao's Syracuse measure
\cite{Tao2022}. Each claim was read in full, its central argument
reconstructed independently, and checked computationally wherever it
admits a numerical test; several of the more delicate judgment calls
were also cross-checked against an independent large-language-model
mathematical reviewer. None of the twelve survives scrutiny as a valid
proof or disproof.

The point of this note is not the bare fact that twelve amateur claims
are wrong --- that would be unremarkable and uninteresting on its own.
The point is \emph{how} they fail. Section~\ref{sec:taxonomy} groups
the twelve cases into five categories by the shape of their central
error, together with a brief summary of each case; three categories,
covering ten cases, converge on the same structural obstruction from
routes that look, on the surface, entirely unrelated to one another.
Section~\ref{sec:discussion} discusses what this convergence suggests
about why the conjecture resists these particular styles of attack,
and what the two exceptions to the pattern reveal by contrast.
Section~\ref{sec:conclusion} concludes.

To be clear about what this note is \emph{not}: it is not a claim
about the authors' competence or good faith, several of whom are
explicit and careful about the limits of their own arguments elsewhere
in the same text (correctly labeling an open step a ``conjecture''
rather than a ``theorem'', for instance). The focus throughout is
strictly the mathematical content of the specific error identified in
each case.

\section{Methodology}
\label{sec:methodology}

Each of the twelve claims was: (i) read in full, not just abstract and
conclusion; (ii) reconstructed as an explicit chain of premises and
inferences, so that the exact step carrying the logical weight of
``proves'' or ``disproves'' could be isolated; (iii) checked
computationally whenever the claim, or the step suspected to be the
gap, admits a numerical test --- either by direct reimplementation of
the claimed identity, or by constructing a targeted counterexample; and
(iv), in the handful of cases where the mathematical judgment involved
was subtle enough to warrant a second opinion, cross-checked against an
independent large-language-model reviewer with no access to our own
working conclusion in advance. Full technical detail for every case ---
the exact reconstructed argument, the computational check, and where
applicable the source code --- is kept in individual case files
referenced in Section~\ref{sec:taxonomy}; this note summarizes rather
than reproduces that material.

\section{Taxonomy of recurring error patterns}
\label{sec:taxonomy}

An initial pass at this taxonomy, built only from each case's title and
a one-line status summary, proposed seven categories. It did not
survive contact with the full texts: one case (Boyle, \S\ref{sec:catC})
fit none of the seven, and reading the remaining cases in full showed
that it, together with four others, was in fact the same error dressed
in unrelated formal apparatus. The taxonomy below reflects the full
reading and has five categories, one of which --- Section
\ref{sec:catC} --- accounts for five of the twelve cases by itself.

\subsection{Category A: circular reasoning}
\label{sec:catA}

The argument's central object, or its central step, is well-defined or
valid only under the assumption that the conjecture already holds ---
so the ``proof'' does not establish termination, it merely restates it
in different notation. Three cases:

\textbf{Getachew (2025)} \cite{Getachew2025} constructs a reverse tree
of the Collatz map and defines a ``parent'' relation that turns out to
be, term for term, identical to the forward Collatz map itself
(\(\mathrm{parent}(m) = m/2\) for even \(m\), \(3m+1\) for odd \(m\)).
The paper's Lemma~4.3, ``every path back to the root is finite,'' is
then logically equivalent to the conjecture itself, not a consequence
of the tree's acyclicity and unique-parent structure as claimed:
acyclicity and a unique parent guarantee that a path never repeats or
branches, never that it is finite. The paper's covering theorem
(Theorem~5.1) fails for a parallel reason: it proves a universal
arithmetic decomposition (odd number times a power of two) that we
confirmed computationally never actually uses the branching rule.

\textbf{Spencer (2025a)} \cite{Spencer2025a} builds a different reverse
tree, indexed by residue classes modulo \(2\cdot 3^J\), with a
ternary-counter argument partially formalized in Lean. It correctly
proves that every primitive residue class is ``occupied'' at every
finite scale \(J\) --- but its Theorem~15.3 then concludes that every
specific integer has a finite reverse address, a step never justified:
``residue occupied'' does not imply ``exact value reached.'' We
confirmed this gap directly: the residue class of \(n=27\) is occupied
in the real tree built to depths 6 and 8, but only by representatives
on the order of \(10^6\)--\(10^8\), never by 27 itself.

\textbf{Yun (2026)} \cite{Yun2026} contains three independent circular
arguments, any one of which alone invalidates the paper's claim. Its
central ``rank function'' \(r(x) := \min\{n \in \NN : f^{(n)}(x) = 1\}\)
is defined only for values of \(x\) whose orbit already reaches 1 ---
if some orbit never reached 1, exactly the possibility the conjecture
needs to exclude, \(r(x)\) would simply not exist, so the well-founded
descent argument built on it presupposes its conclusion. We confirmed
this by applying the identical construction to \(g(x)=2x\), a map that
provably diverges for every \(x>0\): the analogous rank function is
undefined for all 18 tested values, exactly as it would be for \(f\) if
some orbit failed to reach 1. Separately, the paper's cycle-exclusion
argument (Theorem~6.1) asserts that ``the only known fixed point is
\(x=1\)'' as if this excluded other cycles, and its ``informational
compression'' argument (\S5, \S9.2.3) is a non sequitur about
cardinalities of infinite sets that is compatible with any long-run
behavior, not specifically convergence to 1.

\subsection{Category B: illegitimate generalization from a single sequence}
\label{sec:catB}

The argument shows, correctly, that one specific hand-picked sequence
of increasing and decreasing steps has a certain property, then
extends the conclusion to all sequences without bounding how long an
adversarial run of increasing steps could be --- exactly the step the
structural obstruction makes hard. Two cases:

\textbf{Santos (2018)} \cite{Santos2018} introduces a binary
``moves''/``strong reductions'' notation equivalent to the ordinary
2-adic valuation sequence of a Collatz orbit. Its core argument
(\S2.6) picks a hypothetical \(K>30\), \emph{assumes a specific fixed
sequence} of steps (two increasing moves followed by three decreasing
ones, with hand-chosen divisors), shows algebraically that \(K\)
decreases along \emph{that particular sequence}, and then generalizes
this to every \(K>30\) without ever showing that adversarial sequences
of increasing steps cannot be made arbitrarily long.

\textbf{Halemane (2025), the CTUHSK theorem} \cite{Halemane2025}
reformulates the standard reverse tree with new notation
(Binary-Exponential-Ladders, operators \(D^\#/U^\#\)) but adds nothing
structurally new. Its necessary condition is a tautology (the set
\emph{defined} as reaching the trivial cycle does reach the trivial
cycle); its sufficient condition, which carries the real weight of
excluding extra cycles and divergent chains, has a decisive hole: the
reductio argument (\S10.2.1) assumes the predecessor of a minimal
element of a hypothetical cycle is given by the paper's formula with
exponent \(v=1\) specifically, when that formula in fact has infinitely
many valid solutions, one per odd exponent \(v=1,3,5,\dots\); only
\(v=1\) yields a smaller value (violating the assumed minimality),
while every \(v\ge 3\) yields a larger one, generating no contradiction
at all. We confirmed computationally that this is a general algebraic
pattern, not an artifact of the paper's chosen numerical example.

\subsection{Category C: an aggregate or asymptotic property mistaken for a deterministic, pointwise guarantee}
\label{sec:catC}

This category is the paper's central empirical finding. Five
superficially unrelated arguments --- built respectively on Baker's
theorem, the classical random-walk heuristic, a 2-adic Lyapunov bound,
Big-O asymptotic notation, and static 2-adic density --- turn out to
commit the identical error: a property that holds only as an average,
an ensemble expectation, or an asymptotic bound is substituted, without
justification, for a guarantee that must hold deterministically, at
every point, along one specific trajectory. Three of the five papers
below cross-reference this exact pattern in each other, independently
of this catalog, which suggests it is a natural attractor for
Collatz-proof attempts rather than an artifact of our classification.

\textbf{Mohammed} \cite{Mohammed2024} builds a correct geometric density
sieve and correctly excludes 1-cycles, then attempts to exclude
nontrivial cycles via Baker's theorem on linear forms in logarithms.
The cycle equation genuinely forces \(M/P \to \log_2 3\) for a real,
self-consistent cycle with large elements --- but the paper substitutes
\(M \approx 2P\), the \emph{ergodic, average} expectation for a generic
odd integer (already known in this project as a baseline fact, not a
cycle constraint), producing a check that constrains nothing.

\textbf{Boyle (2026)} \cite{Boyle2026} derives, via the standard
random-walk heuristic, that the ensemble-average fraction of even steps
along Collatz orbits is \(2/3\) --- valid, at best, as an average over
many random starting values \(n\) --- and then substitutes this
constant directly into an equation (Lemma~4.2) about one specific,
fixed, hypothetical divergent trajectory, deriving a ``contradiction''
that depends entirely on this illegitimate substitution. We confirmed
computationally that even among orbits known to converge, the fraction
of even steps realized by any one individual trajectory varies
substantially around \(2/3\) (mean 0.677, standard deviation 0.031 over
3000 sampled trajectories) --- it is an aggregate statistic, not a
constraint on any single orbit.

\textbf{Tynski (2026)} \cite{Tynski2026} builds an extensive and mostly
correct formal apparatus (Clifford algebra, an explicit trace,
established cycle-exclusion identities) around a ``master theorem''
unifying the Riemann hypothesis and the Collatz conjecture, but the
exclusion of divergent orbits rests entirely on an axiom asserting a
\emph{deterministic} Lyapunov bound
\(\log_2 n_m \le \log_2 n_0 - m\delta + C\sqrt m\) for a universal
constant \(C\), justified circularly (the text derives the needed
decay ``once \(A_m/m > \log_2 3\) is forced,'' the very property being
argued for) and supported only by numerical verification up to
\(n_0 \le 3\times 10^4\). We tested the claim directly beyond that
range: the minimal \(C\) required already exceeds the claimed universal
bound of 5 for seeds \(10^4 \le n_0 < 10^9\) (worst case 5.08), and
known long-delay record-holders require \(C\) up to 8.34 --- a direct
computational refutation, not merely an absence of proof.

\textbf{The Syzdykov dispute (2025--2026)} \cite{Syzdykov2025,Syzdykov2026,LafontaineCheong2026}
is a three-paper exchange rather than a single claim. The original note
\cite{Syzdykov2025} asserts \(O(f(n)) \le n\) for the Collatz function
\(f\), substitutes \(n = 2^m-1\), and declares
\(O(f(2^m-1)) \approx 3^m > 2^m - 1\) a contradiction refuting the
conjecture. Lafontaine \& Cheong \cite{LafontaineCheong2026} correctly
identify the category error: \(O(\cdot)\) denotes a set of functions,
not a scalar, so \(O(f(n)) \le n\) is not well-formed; and even under a
charitable pointwise reading, an asymptotic bound (``for some \(C,n_0\),
for all \(n \ge n_0\)'') licenses no conclusion about the exact value of
one specific \(n\) --- the same confusion of an existential, aggregate
bound with a pointwise, deterministic one seen in the other four cases
in this category. Syzdykov's reply \cite{Syzdykov2026} does not engage
this point; it redefines \(O(f(n))\) after the fact to mean a scalar
orbit length and reasserts the original conclusion as ``already
proved.''

\textbf{Barghout (2019)} \cite{Barghout2019} correctly proves, by
elementary induction, that the 2-adic valuation of even integers (or of
values produced by \(3n+1\)) follows a fixed, deterministic density
pattern --- a static fact about the \emph{set} of even integers, not
about any specific orbit. Its convergence argument then imports this
static density as if it were the \emph{visit statistic} of one specific
dynamical orbit, implicitly assuming the orbit equidistributes ---
precisely the unproved content of the conjecture, not a consequence of
the density lemmas. We confirmed by direct simulation that individual
orbits (e.g., \(n=9663\)) can grow to roughly 2800 times their starting
value before eventually decreasing, illustrating concretely why an
``on-average'' argument constrains no single trajectory.

\subsection{Category D: an elementary error unrelated to the structural obstruction}
\label{sec:catD}

\textbf{Roif (2026)} \cite{Roif2026} constructs a three-point
topological space \(S=\{0,2,1\}\), reinterprets the integer 2 as an
``empty container'' \(\{\}\), and proves a lemma asserting that any
subset of the positive integers with asymptotic density zero must be
empty. This is false as a general mathematical fact and is the only
error in the catalog with no connection to the conjecture's real
difficulty: density zero does not imply finiteness, let alone
emptiness. The set of perfect squares, or of powers of 2, has density
zero and is plainly infinite. This single elementary error collapses
the paper's entire closing argument (its Theorem~9.1 and
Corollary~9.2, which would otherwise complete the proof), independent
of the remaining formal apparatus (Clifford-adjacent identities, a
Lyapunov-style wave function), which is a correct but by-itself
insufficient restatement of results already standard in the
literature.

\subsection{Category E: a genuine rigor gap without a computational counterexample}
\label{sec:catE}

\textbf{Spencer (2026)} \cite{Spencer2025b}, by the same author as the
Category A case in \S\ref{sec:catA}, uses a different reverse tree
(infinite-degree, via the inverse of the accelerated map). Unlike its
companion paper, building the real tree and testing coverage up to
\(10{,}000\) found \emph{no} gap in practice --- apparent coverage
failures at smaller computational budgets converged exactly to the
root once checked without a magnitude cutoff. We nonetheless identified
a genuine gap in the written argument: Theorem~14.1 proves occupation
of residue \emph{classes}, and Theorem~14.2 concludes occupation of
individual \emph{elements} without justifying the step --- the same
logical shape as the Category A case, but one that, so far, has not
manifested as an actual failure under computational search. This
distinction matters: not every gap in a Collatz proof attempt is
equally fatal, even within the same author's closely related
arguments.

\section{Discussion}
\label{sec:discussion}

Ten of the twelve cases --- all of Categories A, B, and C --- fail, by
routes that look unrelated on the surface, at exactly the same point:
none of them establishes a genuine bound on how long an adversarial run
of increasing steps can be, the structural obstruction that no known
method resolves. Category C alone accounts for five of the ten, and is
the most striking part of the pattern: the same conceptual error ---
substituting an aggregate, average, or asymptotic property for a
pointwise, deterministic guarantee --- survives being dressed in five
unrelated formal vocabularies (elementary number theory via Baker's
theorem, the classical random-walk heuristic, a 2-adic Lyapunov bound,
Big-O asymptotic notation, and static 2-adic density). This is not an
artifact of our own classification: three of the five source documents
already cross-reference each other, and two of them cross-reference the
random-walk heuristic explicitly, as instances of the same confusion,
independently of this catalog.

The two cases outside this pattern are informative by contrast. Roif's
error (Category D) is an isolated, elementary mistake with no bearing
on the conjecture's actual difficulty --- a reminder that not every
error in a Collatz proof attempt is evidence about \emph{why} the
conjecture is hard, some are simply wrong for unrelated reasons.
Spencer's second paper (Category E) shows the opposite lesson: an
identical logical shape to a Category A failure (residue-class
occupation conflated with individual-element occupation) can, in one
instance, fail to manifest as an actual computational counterexample
--- underscoring that the taxonomy classifies argument \emph{structure},
not severity, and that structurally similar gaps are not always equally
fatal.

We note, finally, that the broader collection of Collatz-related papers
surveyed alongside these twelve (outside the scope of this note)
contains many papers that make no claim of a complete proof at all, and
that are careful to label open steps as conjectural. The twelve cases
here were selected precisely because they make the strong claim; the
taxonomy above should not be read as suggesting that claimed-proof
papers are typical of the wider literature on the conjecture.

\section{Conclusion}
\label{sec:conclusion}

All twelve reviewed claims of a complete proof or disproof of the
Collatz conjecture fail under independent scrutiny, and none is close
to valid. The value of this catalog is not the count but the
concentration: half of the twelve cases, despite using entirely
different mathematical vocabularies, commit the identical error of
mistaking an aggregate or asymptotic property for a deterministic,
pointwise guarantee, and three more commit one of two closely related
errors (circularity, illegitimate generalization from a single
sequence) that converge on the same known structural obstruction. This
regularity is itself a small piece of evidence about why the conjecture
resists elementary attack: the failure modes are not diverse, they
repeat.

\section*{Acknowledgments}

Parts of the verification described in this note (case reconstruction,
computational checks, and cross-checking of subtle mathematical
judgment calls) were carried out with the assistance of Claude (models
in the Claude 4 and 5 families, Anthropic), used as a research
assistant under the direction and review of the author.

\begin{thebibliography}{99}

\bibitem{LagariasSurvey} J.\ C.\ Lagarias (ed.), \emph{The Ultimate
  Challenge: The $3x+1$ Problem}, American Mathematical Society, 2010.

\bibitem{Tao2022} T.\ Tao, \emph{Almost all orbits of the Collatz map
  attain almost bounded values}, Forum of Mathematics, Pi \textbf{10}
  (2022), e12.

\bibitem{Getachew2025} E.\ S.\ Getachew, \emph{Unfolding the Collatz
  Tree: An Indirect Structural Proof}, Research in Mathematics
  (Taylor \& Francis), 2025.

\bibitem{Spencer2025a} M.\ E.\ Spencer, \emph{Finite Block Exhaustion
  and Rooted Occupancy in the Inverse Collatz System: A No-Alternative
  Proof via Ternary Counters, Primitive Carry, and Forced Thread
  Displacement}, academia.edu preprint, 2025.

\bibitem{Spencer2025b} M.\ E.\ Spencer, \emph{Rooted Surjectivity from
  the Invariant E/O Refinement System}, academia.edu preprint, 2026.

\bibitem{Santos2018} O.\ de O.\ Santos, \emph{Proving the Collatz
  Conjecture with Binaries Numbers}, Pure and Applied Mathematics
  Journal \textbf{7}(5) (2018), Science Publishing Group.

\bibitem{Halemane2025} K.\ P.\ Halemane, \emph{Collatz-Thwaites-Ulam-Hasse-Syracuse-Kakutani
  (CTUHSK) Theorem: Convergence of Collatz ($3n+1$) Sequence to the
  Trivial Cycle Proved}, engrXiv preprint, 2025.

\bibitem{Mohammed2024} A.\ Mohammed, \emph{Structural Analysis, Dynamic
  Density Sieve, and Logarithmic Contraction of Collatz Sequences},
  preprint, Kafr El-Sheikh University, 2024.

\bibitem{Boyle2026} D.\ G.\ Boyle, \emph{The Collatz Conjecture is
  True}, rxiverse.org preprint, 2026.

\bibitem{Roif2026} A.\ Roif, \emph{On the Convergence of the Collatz
  Function via Galois Fields, Topological Duality, and Lyapunov Wave
  Analysis}, academia.edu preprint, v4, 2026.

\bibitem{Yun2026} Y.\ H.\ Yun, \emph{A Structural Proof of the Collatz
  Conjecture via Non-Repeating Trajectory and Recursive Decay}, OSF
  preprint, Zero Theoretical Physics Lab, 2026.

\bibitem{Tynski2026} K.\ Tynski, \emph{A Common Proof of the Riemann
  Hypothesis and the Collatz Conjecture}, academia.edu preprint, 2026.

\bibitem{Syzdykov2025} M.\ Syzdykov, \emph{Disproof of Collatz
  Conjecture by Using O-notation}, ResearchGate preprint, 2025.

\bibitem{Syzdykov2026} M.\ Syzdykov, \emph{Continued Disproof Sentence
  towards Collatz Conjecture}, ResearchGate preprint, 2026.

\bibitem{LafontaineCheong2026} M.\ Lafontaine and G.\ Cheong, \emph{A
  Note on a Claimed Disproof of the Collatz Conjecture}, ResearchGate
  preprint, 2026.

\bibitem{Barghout2019} K.\ Barghout, \emph{On the Probabilistic Proof
  of the Convergence of the Collatz Conjecture}, Journal of Probability
  and Statistics (Hindawi) \textbf{2019}, Article ID 6814378.

\end{thebibliography}

\end{document}
